%! TEX root = ../main.tex
\documentclass[main]{subfiles}

\begin{document}
\section{結果と考察}
\subsection{CMT}
\begin{table}[H]
 \begin{center}
   \begin{tabular}{|l|r|r|r|}
    \hline
    MT-type                       & PSDS1 & PSDS2 & mpAUC \\
    \hline
    Mean Teacher(MT)            & 0.369 & 0.581 & 0.720 \\
    Confident Mean Teacher(CMT) & 0.297 & 0.60  & 0.689 \\
    $CMT_{neg}$                     & 0.369 & 0.543 & 0.661 \\
    $CMT_{scale}$                 & 0.368 & 0.621 & 0.699 \\
    \hline
  \end{tabular}
   \caption{Mean Teacher比較}
   \label{table:mtAcc}
\end{center}
\end{table}

表\ref{table:mtAcc}に各手法の精度比較を示す．MTと比較して，CMTはPSDS2が向上したが，PSDS1及びmpAUCは低下した．
対して正負例を学習する$CMT_{neg}$は，CMTで見られたPSDS1の低下は改善されたものの，その他の指標は精度が悪化した．
正負例の均衡化処理を行う$CMT_{scale}$は，MTと同等のPSDS1を維持しつつ，PSDS2においてCMTを上回る性能を示した．ただし，mpAUCはMTに及ばなかった．\par

CMTにおけるPSDS1の低下は，非イベント区間の情報（負例）の欠落で時間的な予測が不安定になったためと考えられる．
$CMT_{neg}$では負例の導入によりこの問題を解決したが，正例に対する負例の割合が過多となり，逆にクラス分類（PSDS2など）の予測が不安定化した．
よって，$CMT_{scale}$における正負例の量の調整が，PSDS1，2の安定に有効であると言える．\par

\subsection{クラス別CMT}
\label{classwise}
ここでCMTについて，クラスごとの精度を見る．表\ref{table:classWise}に各クラスにおけるPSDS2及びmpAUCの結果を示す．\par
本稿の将来的な目標である異常検知において，異常発生の重要な文脈となる「流水（Running water）」と「足音（Footsteps）」に着目する．
ベースラインのMTと比較して，提案手法である$CMT_{scale}$ は，流水のPSDS2において$0.88$から$0.91$へ，足音のmpAUCにおいて$0.68$から$0.71$へと，それぞれ精度向上を達成した．\par
また，表\ref{table:classWise}の他手法と比較して{5\%}以上の変化があったクラスについて，向上は青，低下は赤で強調した．向上したクラスは音が長く持続するが，低下したクラスでは瞬間的な音が続く．
このことから，提案手法は長期的な音が捉えられる代わりに，瞬間的な音を取りこぼすと考えられる．

\begin{table}[H]
  \centering
  \footnotesize
  \setlength{\tabcolsep}{3pt}
  \begin{tabular}{|l|c|c|c||l|c|c|c|}
    \hline
    \multicolumn{4}{|c||}{PSDS（DESED）} & \multicolumn{4}{c|}{mpAUC（MAESTRO）} \\
    \hline
    class & MT & CMT & $CMT_{s}$ & class & MT & CMT & $CMT_{s}$ \\
    \hline
    Alarm\_bell & .868 & .937 & \textbf{.947} & Birds\_singing & .971 & .969 & \textbf{.972} \\
    \textcolor{RoyalBlue}{\textbf{Blender}} & .870 & .885 & \textbf{.920} & Brakes\_squeak & .416 & .398 & \textbf{.463} \\
    Cat & .881 & \textbf{.929} & .911 & \textcolor{BrickRed}{\textbf{Car}} & .825 & \textbf{.866} & .817 \\
    \textcolor{BrickRed}{\textbf{Dishes}} & .501 & \textbf{.641} & .602 & Children voices & \textbf{.916} & .868 & .881 \\
    Dog & .848 & \textbf{.893} & .883 & Cutlery/dishes & .595 & .584 & \textbf{.617} \\
    Elec\_shaver & .965 & \textbf{.989} & .984 & Footsteps & .681 & .696 & \textbf{.710} \\
    Frying & .958 & .948 & \textbf{.970} & Large\_vehicle & .653 & .659 & \textbf{.681} \\
    Running\_water & .887 & .902 & \textbf{.912} & \textcolor{BrickRed}{\textbf{Metro approach}} & \textbf{.833} & .807 & .782 \\
    Speech & .775 & .853 & \textbf{.862} & Metro leaving & \textbf{.768} & .756 & .740 \\
    Vacuum\_clean & .901 & .912 & \textbf{.948} & People talking & \textbf{.665} & .662 & .662 \\
    & & & & \textcolor{RoyalBlue}{\textbf{Wind\_blowing}} & .599 & .324 & \textbf{.744} \\
    \hline
  \end{tabular}
  \caption{クラス別精度比較（$CMT_{s}$は$CMT_{scale}$の略）}
  \label{table:classWise}
\end{table}

\subsection{MixStyle}
\begin{table}[!h]
 \begin{center}
   \begin{tabular}{|l|r|r|r|}
    \hline
    CNN        & PSDS1 & PSDS2 & mpAUC \\
    \hline
    -          & 0.375 & 0.581 & 0.709 \\
    \checkmark & 0.373 & 0.596 & 0.745 \\
    \hline
  \end{tabular}
  \caption{MixStyle+CNNの有無}
  \label{table:mixstyle}
\end{center}
\end{table}
表\ref{table:mixstyle}より，CNNでの重み付けによってPSDS2とmpAUCが向上した．なお，半教師あり学習手法としてMTを用いた．
 これはドメイン汎化がより促されたか，モデルの入力が変化し，一貫性の制約が強まったためと考えられる．
 ただ，重みづけは実際には働いていない可能性がある．また，従来のMixStyleは混合した特徴量のみを用いるが，提案は残差接続を行う．これが摂動として働き，一貫性の制約を強めた可能性もある． \par

 \subsection{CMT+MixStyle}
 \begin{table}[!h]
 \begin{center}
   \begin{tabular}{|l|r|r|r|r|r|}
    \hline
    MT-type & PSDS1 & PSDS2 & mpAUC \\
    \hline
    MT   & 0.382 & 0.607 & 0.704 \\
    \hline
    CMT & 0.387 & 0.637 & 0.722 \\
    \hline
    $CMT_{scale}$  & 0.393 & 0.632 & 0.727 \\
    \hline
  \end{tabular}
  \caption{CMTとMixStyleを併用した結果}
  \label{table:mergeAcc}
\end{center}
\end{table}

表\ref{table:mergeAcc}に，各種MTをMixStyleと併用した場合の結果を示す． 
MTとCMTはPSDS1，2が向上した一方で，mpAUCは低下した．
$CMT_{scale}$はMTに比べて精度が向上したが，CMTとの差は僅かである．
これはMixStyleの精度への影響がCMTより強いためだと考えられる．\par 

表\ref{table:mtAcc}と比較する．MixStyleの適用により，DESED（合成音・実録音）のみならず，性質の異なるMAESTRO（実録音）を含む精度向上が期待されたが，
MTではmpAUCが低下した． 
PSDS1，2の向上は，MixStyleが特徴空間に摂動を与え，教師及び生徒モデル間の一貫性を強めたためと考えられる．\par

\subsection{まとめ}
一連の結果より，提案手法であるCMTとMixStyleの改良を統合した場合，MTと比較して全ての指標を向上させたが，CMTとの差はごく僅かだった．
正例に加え負例を活用する$CMT_{scale}$の有効性は示されたように思うが，あくまでMTの範疇である．
また，MixStyleと軽量CNNによる重み付けも，精度向上に一定の寄与をすることが確認された．\par 

 \end{document}